\section{相关工作}\label{sec:related}

\subsection{自我改进与协同演化}\label{sec:related:treesearch}

reflection 与 self-refinement 把模型生成的 critique 反馈给后续尝试
\cite{madaan2023self,shinn2023reflexion}。它们改进 artifact 或 trajectory,
但 evaluator 与权限结构大体固定。Constitutional ARI-RQGM 则把 prompt
role 与 utility policy 视为 versioned institutional artifact,其替换必须
具备证据、角色分离与合法 transition。

最接近的概念祖先是 Red Queen G\"odel Machine
\cite{iacob2026red},其中 agent 与 evaluator 协同演化,scoring function
也可以改变。ARI-RQGM 采用 co-evolution 与 adversarial replay,并增加固定
宪制、epoch transaction、single-writer registry、target-bound sanction、
clean-room replacement,以及 mutable policy 与 immutable verification 的
明确分离。本文不把原 RQGM 已有的协同演化和边界更新作为新颖点。概念上
的差异是用权限矩阵和完整组件生命周期表围住可变制度,再在同一来源图上连接退役
修复与固定证据检查。系统贡献在于这一组合及其可检查边界,而非单个机制。

Constitutional AI 用书面原则进行自我批评和人工智能反馈
\cite{bai2022constitutional}。本文的“宪法”在性质上不同:它固定模型输出
不能豁免的权限矩阵、状态转换和证据格式。AI Control 在有意规避监督的
前提下评估可信与不可信模型分离的协议 \cite{greenblatt2023ai}。本文范围
更窄,尚未做定量的共谋控制评估;但写作、审稿、裁决与注册表变更不会由
同一权限掌握。

\subsection{自主研究系统}\label{sec:related:autonomous}

AI Scientist 展示了从 ideation 到 experimentation 与 manuscript
preparation 的 end-to-end loop \cite{lu2024ai};后续工作分析了走向自动
科学发现的更广趋势 \cite{lu2026towards}。AIDE 把实验 code 改进视为反复
agent 问题 \cite{jiang2025aide}。这些系统确立了 persistent artifact、
executable evaluation 与自动写作的价值。它们主要优化 research artifact;
Constitutional ARI-RQGM 的对象则是提出、评估和撰写这些 artifact 的制度。

Agent Laboratory 研究作为 research assistant 协作的 model role
\cite{schmidgall2025agent},MLE-bench 为 machine-learning engineering
agent 提供 executable task \cite{chan2024mle}。constitutional boundary
不依赖底层 actor 使用 prompted reasoning \cite{wei2022chain} 还是
large-scale candidate generation \cite{li2022competition};它治理的是
actor version、authority 与 evidence,而不是内部生成过程。

Story2Proposal 主张用契约把 narrative 科学主张连接到 executable proposal
与 evidence \cite{qian2026story2proposal}。governed paper archive 把同一
一般洞见用作 immutable writer anchor:演化中的 reviewer 不能让违背执行
证据的陈述成为事实。

\subsection{评估与证据}\label{sec:related:eval}

PaperBench 通过结构化 rubric 评估 agent 能否复现研究论文
\cite{starace2025paperbench}。对治理而言,其核心启示是评估主张需要明确的
task structure、可观测 artifact 与 grading boundary。Constitutional
ARI-RQGM 把这一点从最终 outcome 扩展到制度史:attack、defense、judgment、
transition 与 repair 都成为保留 provenance 的 typed record。

adversarial evaluation 常用于强化固定 model 或 policy。在这里,validated
attack 还有两项功能:为 serving component 的问责提供证据,并成为 candidate
successor 的 replay case。因此,同一个故障既是关于 incumbent 的证据,
也是替代者的 entrance test。

\subsection{本工作的定位}\label{sec:related:foundations}

最小权限、完全仲裁和权限分离是经典的保护设计原则
\cite{saltzer1975protection}。确定性状态机为从相同已提交输入序列恢复
相同状态奠定基础 \cite{schneider1990implementing}。近期事件溯源智能体
设计同样采用只追加事件和确定性投影 \cite{nakajima2026the}。因此,本文
不把单写者状态、重放或只追加历史本身视为发明。技术任务是把它们一致地
连接到 RQGM 纪元、责任归属、退役修复和稿件证据检查。

\href{https://www.w3.org/TR/prov-o/}{W3C PROV-O 建议书}规定可互操作的
实体--活动--主体来源模型。
\href{https://www.usenix.org/conference/usenixsecurity19/presentation/torres-arias}{in-toto}
使用签名的软件供应链步骤证据，
\href{https://docs.sigstore.dev/logging/overview/}{Sigstore Rekor} 提供可由
第三方审计的仅追加透明日志。它们具有当前本地未签名检查点所没有的来源
互操作性或密码学/见证完整性；反之，它们不定义本文的 RQGM 退役闭包与
效用重评分。

\begin{table*}[t]
  \centering
  \scriptsize
  \begin{tabularx}{\textwidth}{@{}p{0.12\textwidth}*{7}{X}@{}}
    \toprule
    \textbf{系统谱系} & \textbf{可变对象} & \textbf{固定基础} &
    \textbf{角色权限} & \textbf{原子状态} &
    \textbf{历史完整性} & \textbf{修复} & \textbf{威胁/证据} \\
    \midrule
    RQGM \cite{iacob2026red} & 主体、评估者、效用 &
      纪元内评估 & 非重点 & 纪元更新 &
      非重点 & 协同演化选择 &
      概念/实验协同演化 \\
    事件溯源主体 \cite{nakajima2026the} & 图行为与状态 &
      事件契约 & 意图/状态分离 & 确定性投影 &
      只追加来源 & 分叉/重放 &
      可审计架构 \\
    Story2Proposal \cite{qian2026story2proposal} & 科学叙事与提案 &
      结构化契约 & 非重点 & 契约阶段 &
      证据链接 & 契约式修订 &
      写作支架评估 \\
    AI Control \cite{greenblatt2023ai} & 控制协议 &
      可信监视/协议 & 可信/不可信分离 &
      协议特定 & 监视记录 & 协议响应 &
      有意颠覆与实证控制 \\
    来源/透明系统 & 实体、活动、签名陈述 &
      来源模式、密钥、日志协议 & 认证生产者/验证者 &
      签名步骤或追加操作 & 可互操作或有见证的历史 &
      验证而非语义退役修复 & 密码学完整性 \\
    本文 & 提示词、组件、效用策略 &
      固定内核与验证器 & 进程内权限矩阵 &
      边界事务 & 语义事件摘要；无外部锚 &
      退役闭包与重评分 &
      五个样例及回归试验；无共谋研究 \\
    \bottomrule
  \end{tabularx}
  \caption{按组件定位。单元格有内容不表示目标不同的系统可直接比较。
    本文的新颖性主张是集成，而非发明各个机制。}
  \label{tab:related-comparison}
\end{table*}

透明日志与软件供应链来源询问哪些字节和身份被提交；引用监视器与能力
安全询问哪个主体能对哪个对象执行何种操作；安全自动机与状态机方法询问
所有执行轨迹是否服从安全语言 \cite{schneider2000enforceable}。这些是相关谱系，并非本文的新颖性主张。
当前实现采用最小权限与完全仲裁 \cite{saltzer1975protection} 以及确定性
状态投影 \cite{schneider1990implementing}，但既没有独立见证的日志头，
也没有 OS 强制的角色能力；其复合纪元、事务、配额和修复状态也未经过
模型检查。

Constitutional ARI-RQGM 既不宣称自主科学优越,也不是自修改智能体的一般
解法。它解决一个更窄的系统问题:如何让 model-defined research role 与
utility policy 在运行中改变,同时不赋予这些角色批准自身变更的权限?
其回答结合 co-evolution、deterministic constitutional enforcement、
transactional state、provenance-aware repair 与固定 evidence gate。贡献
在于这些机制的组合,以及明确区分其保证与仍需实证研究的主张。
